\documentclass[11pt,a4paper]{amsart}
\pdfinfoomitdate=1
\pdftrailerid{}
\pdfsuppressptexinfo=15
\usepackage[T1]{fontenc}
\usepackage{lmodern}
\usepackage[margin=27mm]{geometry}
\usepackage{microtype}
\usepackage{amsmath,amssymb,mathtools}
\usepackage{booktabs}
\usepackage{xcolor}
\usepackage[numbers,sort&compress]{natbib}
\usepackage[colorlinks=true,linkcolor=blue!55!black,citecolor=blue!55!black,urlcolor=blue!55!black]{hyperref}
\usepackage[nameinlink,noabbrev]{cleveref}
\raggedbottom

\newtheorem{theorem}{Theorem}[section]
\newtheorem{proposition}[theorem]{Proposition}
\newtheorem{corollary}[theorem]{Corollary}
\newtheorem{lemma}[theorem]{Lemma}
\theoremstyle{definition}
\newtheorem{definition}[theorem]{Definition}
\theoremstyle{remark}
\newtheorem{remark}[theorem]{Remark}

\newcommand{\F}{\operatorname{Fix}}
\newcommand{\Per}{\operatorname{Per}}
\newcommand{\Z}{\mathbb Z}
\newcommand{\E}{\mathbb E}
\newcommand{\Var}{\operatorname{Var}}
\newcommand{\one}{\mathbf 1}

\title[Least-valuation digit erasure]{Least-Valuation Digit Erasure:\\
Exact Transient Profiles Beyond Periodic Data}
\author{Anonymous}
\date{Internal preprint, 29 August 2026}
\subjclass[2020]{37P35, 11A63, 60F05, 05A15}
\keywords{finite dynamical system, p-adic digit sum, absorption time, transient profile, local limit theorem, dynamical zeta}

\begin{document}

\begin{abstract}
For a prime $p$ and $r\geq1$, consider the absorbing map on
$\Z/p^r\Z$ that subtracts the place value of the least significant nonzero
base-$p$ digit.  We give an exact digit-vector normal form: every step lowers
the digit sum by one, so the hitting time of zero is precisely the base-$p$
digit sum.  Consequently the complete transient-depth polynomial is
\[
 (1+u+\cdots+u^{p-1})^r.
\]
We derive every depth-layer cardinality, the sharp global absorption depth,
symmetry and unimodality, exact moments, and the fixed-base central and local
limit laws.  All members of the family have the same periodic-point sequence
and Artin--Mazur zeta, whereas the transient profile recovers both $p$ and
$r$.  This gives a small exact model in which periodic data are completely
blind but finite-time basin geometry is rigid.  Classical digit-sum and
local-limit results are treated as owned background; the residual result is
their coupling to this explicit finite dynamical system.
For $p=2$ the update is Wegner's classical rightmost-one clearing step, and
that specialization is expressly not claimed here.
\end{abstract}

\maketitle

\section{Introduction and claim boundary}

Periodic-orbit data are among the most effective invariants of finite and
symbolic dynamical systems.  They can nevertheless discard all transient
information.  We exhibit an elementary family where this loss is total:
every parameter choice has one fixed point and no other recurrent point, but
the complete distribution of entrance times is explicit and determines the
parameters.

The statistic that appears is the ordinary sum of base-$p$ digits.  Its
distribution, asymptotics, and refinements have a substantial independent
literature; for example, Drmota and Gajdosik study general distributional
questions for the sum-of-digits function \cite{DrmotaGajdosik1998}.
Likewise, the central and lattice local limit theorems used below are
classical results for sums of independent variables \cite{Petrov1975}.
There is also a direct algorithmic owner at the binary endpoint:
\[
 E_{2,r}(x)=x-2^{v_2(x)}=x\mathbin{\&}(x-1).
\]
Wegner introduced this rightmost-one clearing step to count binary ones
\cite{Wegner1960}; its iteration count is the binary digit sum.  We do not
claim that specialization or identity.  Our residual contribution is the
general-prime standard-representative dynamics, its full transient
polynomial and parameter recovery, and the contrast between universally
blind periodic data and a rigid transient profile.  A bounded search did
not locate that combined general-$p$ dynamical package; this is not a
priority certification.

There is also a deliberate internal distinction from earlier work on
digit-weight automatic orbit closures.  Here the phase space is the finite
local ring itself, time is an absorbing arithmetic update, and the invariant
is the entire hitting-time profile.  No countable subshift, cellular
endomorphism monoid, or Cantor--Bendixson analysis is used.

\section{The erasure map and its digit normal form}

Fix a prime $p$ and an integer $r\geq1$.  We use the standard representatives
$0,1,\ldots,p^r-1$ of $R_{p,r}=\Z/p^r\Z$.  For nonzero $x$, let $v_p(x)$ be
the largest $j$ for which $p^j$ divides $x$.  Thus the choice of standard
representative is part of the definition; no representative-independent ring
operation is being asserted.

\begin{definition}
The \emph{least-valuation digit erasure map} is
\[
 E_{p,r}(0)=0,
 \qquad
 E_{p,r}(x)=x-p^{v_p(x)}\quad (x\ne0).
\]
Its absorption time is
\[
 \tau_{p,r}(x)=\min\{t\geq0:E_{p,r}^t(x)=0\}.
\]
\end{definition}

Write the unique base-$p$ expansion
\[
 x=\sum_{j=0}^{r-1}a_jp^j,
 \qquad 0\leq a_j\leq p-1,
\]
and put $s_p(x)=\sum_ja_j$.

\begin{lemma}[digit normal form]\label{lem:normal}
If $x\ne0$ and $j=\min\{i:a_i>0\}$, then $E_{p,r}$ replaces $a_j$ by
$a_j-1$ and leaves every other digit unchanged.  In particular,
\[
 s_p(E_{p,r}x)=s_p(x)-1.
\]
\end{lemma}

\begin{proof}
The digits below $j$ vanish, so $v_p(x)=j$.  Subtracting $p^j$ decreases the
$j$th digit by one and produces neither a borrow nor a carry.  The assertion
about the digit sum follows immediately.
\end{proof}

The choice of the \emph{least} nonzero digit is essential.  It makes the
map triangular in digit coordinates; subtracting the highest nonzero place
would generally create a different ordered traversal, while subtracting an
arbitrary $p$-power can require a state-dependent convention.

\begin{theorem}[complete hitting-time normal form]\label{thm:hitting}
For every $x\in R_{p,r}$,
\[
 \tau_{p,r}(x)=s_p(x),
 \qquad
 \tau_{p,r}(E_{p,r}^t x)=\max\{s_p(x)-t,0\}.
\]
The unique recurrent point is $0$, and the sharp global absorption depth is
\[
 D_{p,r}=(p-1)r.
\]
The unique state at this maximum depth is $p^r-1$.
\end{theorem}

\begin{proof}
By \cref{lem:normal}, every nonzero iterate lowers the nonnegative integer
$s_p$ by exactly one.  It therefore takes exactly $s_p(x)$ steps to reach
digit sum zero, which is equivalent to reaching $0$.  This also proves the
iterate formula and rules out nontrivial recurrence.  The digit sum is at
most $(p-1)r$, with equality exactly when every digit is $p-1$.
\end{proof}

\section{The full transient profile}

Let
\[
 N_{p,r}(k)=\#\{x\in R_{p,r}:\tau_{p,r}(x)=k\}
\]
and call
\[
 H_{p,r}(u)=\sum_{k\geq0}N_{p,r}(k)u^k
\]
the transient-depth polynomial.

\begin{theorem}[depth polynomial and exact layers]\label{thm:profile}
One has
\[
 \boxed{H_{p,r}(u)=(1+u+\cdots+u^{p-1})^r.}
\]
For $0\leq k\leq(p-1)r$,
\[
 \boxed{
 N_{p,r}(k)=
 \sum_{j=0}^{\lfloor k/p\rfloor}
 (-1)^j\binom rj\binom{k-pj+r-1}{r-1},}
\]
where a binomial coefficient with an inadmissible lower argument is zero.
The sequence $k\mapsto N_{p,r}(k)$ is symmetric and unimodal.
\end{theorem}

\begin{proof}
The bijection from ring elements to their $r$ digits, together with
\cref{thm:hitting}, gives
\[
 H_{p,r}(u)=
 \prod_{j=0}^{r-1}\sum_{a=0}^{p-1}u^a.
\]
Using $(1+\cdots+u^{p-1})^r=(1-u^p)^r(1-u)^{-r}$ and extracting the
$u^k$ coefficient gives the displayed inclusion--exclusion formula.

Symmetry follows by complementing every digit, $a_j\mapsto p-1-a_j$.
For unimodality, induct on $r$, the case $r=1$ being immediate.  If $c_r(k)$ denotes the coefficient at
level $r$, extended by zero outside $0\leq k\leq(p-1)r$, then
\[
 c_{r+1}(k)-c_{r+1}(k-1)=c_r(k)-c_r(k-p).
\]
Put $C_r=(p-1)r/2$.  If the integer $k\leq C_{r+1}$, then
$k\leq C_r+p/2$, and hence
$|k-C_r|\leq|k-p-C_r|$.  Symmetry and the inductive monotonicity toward
$C_r$ therefore give $c_r(k)\geq c_r(k-p)$, so the displayed difference is
nonnegative.  Symmetry supplies the nonpositive differences after the
centre and also covers an integral or half-integral centre.
\end{proof}

The polynomial contains more information than merely the maximal depth.  In
fact, it is a complete parameter fingerprint.

\begin{corollary}[transient rigidity]\label{cor:rigidity}
If $H_{p,r}=H_{q,s}$ for primes $p,q$ and positive integers $r,s$, then
$(p,r)=(q,s)$.  More explicitly,
\[
 r=[u]H_{p,r}(u),
 \qquad
 p=1+\frac{\deg H_{p,r}}{[u]H_{p,r}(u)}.
\]
\end{corollary}

\begin{proof}
The coefficient of $u$ counts digit vectors with one unit digit and is
$r$.  The degree is $(p-1)r$, which then recovers $p$.
\end{proof}

\section{Exact laws and asymptotics from a uniform state}

Let $X_{p,r}$ be uniform on $R_{p,r}$.  Its digits are independent and
uniform on $\{0,\ldots,p-1\}$.  Thus the dynamical hitting time is exactly a
sum of iid bounded lattice variables.

\begin{proposition}[moment transform and moments]\label{prop:moments}
For real $\theta$,
\[
 \E e^{\theta\tau_{p,r}(X_{p,r})}
 =\left(\frac{1+e^\theta+\cdots+e^{(p-1)\theta}}p\right)^r.
\]
Moreover,
\[
 \boxed{\E\tau_{p,r}=\frac{r(p-1)}2,\qquad
 \Var(\tau_{p,r})=\frac{r(p^2-1)}{12}.}
\]
\end{proposition}

\begin{proof}
Apply \cref{thm:hitting} and multiply the moment generating functions of the
$r$ independent uniform digits.  A single such digit has mean $(p-1)/2$ and
variance $(p^2-1)/12$.
\end{proof}

Put $\mu_{p,r}=r(p-1)/2$ and
$\sigma_{p,r}^2=r(p^2-1)/12$.

\begin{theorem}[central and local limits]\label{thm:limits}
Fix $p$ and let $r\to\infty$.  Then
\[
 \frac{\tau_{p,r}(X_{p,r})-\mu_{p,r}}{\sigma_{p,r}}
 \Longrightarrow \mathcal N(0,1).
\]
In addition, the span-one lattice local limit holds uniformly in
$k\in\Z$:
\[
 \sup_{k\in\Z}\left|
 \sigma_{p,r}\frac{N_{p,r}(k)}{p^r}
 -\frac1{\sqrt{2\pi}}
 \exp\!\left(-\frac{(k-\mu_{p,r})^2}{2\sigma_{p,r}^2}\right)
 \right|\longrightarrow0.
\]
\end{theorem}

\begin{proof}
By \cref{thm:hitting}, the centered hitting time is a sum of iid bounded
nondegenerate variables, so the classical central limit theorem applies.
The digit distribution has maximal lattice span one and finite moments of
all orders.  The standard iid lattice local limit theorem therefore gives
the uniform statement in the displayed normalization
\cite[Chapter~VII]{Petrov1975}.  The contribution specific to the present
system is the exact identification of that sum with the hitting time.
\end{proof}

\begin{remark}
The fixed-$p$ hypothesis in \cref{thm:limits} is part of the statement.
Joint regimes in which $p$ varies with $r$ require a triangular-array
formulation and are not claimed here.
\end{remark}

\section{Periodic blindness}

For a self-map $F$ of a finite set, write
$\F_n(F)=\#\{x:F^n(x)=x\}$.  The Artin--Mazur zeta is
\cite{ArtinMazur1965}
\[
 \zeta_F(z)=\exp\left(\sum_{n\geq1}\frac{\F_n(F)}n z^n\right).
\]
We use this as an identity in $\mathbb Q[[z]]$.

\begin{theorem}[universal periodic data, rigid transient data]
For every prime $p$, every $r\geq1$, and every $n\geq1$,
\[
 \F_n(E_{p,r})=1,
 \qquad
 \zeta_{E_{p,r}}(z)=\frac1{1-z}.
\]
Thus periodic counts and zeta distinguish no two members of the family,
while the transient-depth polynomial distinguishes all of them by
\cref{cor:rigidity}.
\end{theorem}

\begin{proof}
Theorem~\ref{thm:hitting} shows that every nonzero state strictly decreases
its digit sum until it reaches zero.  Hence only zero can be fixed by an
iterate.  Substitution of $\F_n=1$ into the defining exponential gives
$\exp(\sum_{n\ge1}z^n/n)=(1-z)^{-1}$.
\end{proof}

This contrast is the main dynamical point of the example.  Adding a restart
edge at zero would manufacture a nontrivial cycle but obscure the canonical
absorbing map, so we do not do so.

\section{Exact controls and limitations}

The accompanying deterministic program follows every arithmetic orbit in
five registered parameter lanes and builds the resulting depth histogram.
It cross-checks that exhaustive histogram against an independently convolved
coefficient vector and the inclusion--exclusion formula, then verifies
symmetry, unimodality, exact rational moments, nilpotency depth, fixed data,
and parameter recovery.  The frozen run contains $46{,}319{,}420$ exact
assertions.  These finite checks guard conventions and endpoints; the proofs
above establish the infinite family.

The present result is deliberately narrow.  It does not assert a new theorem
about the classical sum-of-digits distribution or its local limits, does not
classify arbitrary digit-decrement rules, and does not infer global novelty
from a bounded literature search.  In particular it does not re-claim
Wegner's binary clearing step or popcount identity.  Public release and
priority language remain on hold pending specialist review.

\bibliographystyle{plainnat}
\bibliography{references}

\end{document}
